\section{Résumé exécutif}

ENTRADE.MA est un assistant conversationnel bilingue (arabe et français) destiné à orienter les citoyens vers les centres, services et programmes de l'Entraide Nationale du Maroc. Le système repose sur une architecture GraphRAG (Retrieval-Augmented Generation appuyée sur un graphe de connaissances) : les réponses ne sont pas produites librement par un grand modèle de langage, mais construites à partir de données réellement présentes dans deux bases -- un graphe Neo4j qui modélise les relations métier (régions, délégations, communes, centres, services, populations cibles, programmes, institutions, foire aux questions) et une base vectorielle Qdrant qui permet une recherche sémantique hybride sur le même corpus. Le modèle de langage, Qwen2.5-7B-Instruct, est exécuté localement via llama.cpp sur un serveur CPU dédié, sans dépendance à une API externe -- un choix structurant pour la souveraineté des données et le coût d'exploitation, mais qui impose en retour des contraintes de latence fortes qui ont guidé une bonne partie des décisions d'ingénierie décrites dans ce rapport.

Le projet s'est déroulé en deux temps. Une première phase, documentée dans le \emph{Guide de Déploiement} de référence, a posé l'infrastructure : provisionnement des deux machines virtuelles, installation de la couche d'inférence, installation de la couche de stockage, mise en place de l'API d'orchestration, sécurisation réseau, frontend et observabilité. La seconde phase, objet principal du présent rapport, a consisté à faire passer le système d'un prototype fonctionnel à un assistant réellement fiable : ingestion complète et vérifiée des données métier, audit exhaustif de la couverture des données, refonte du pipeline de récupération et de génération pour qu'il réponde \emph{prioritairement} à partir de la base de connaissances tout en sachant reconnaître honnêtement les questions hors domaine, correction méthodique d'une série de défauts réels découverts par des tests de bout en bout (mise en cache inopérante, filtre de pertinence structurellement défaillant, couverture Centre-Service incomplète pour deux catégories de population, ancrage géographique insuffisant, fuite de caractères parasites du modèle, contrôle de langue trop permissif), puis ajout d'une mémoire conversationnelle pour que les questions de suivi elliptiques restent comprises dans leur contexte réel.

Chaque intervention décrite dans ce document a suivi la même discipline : reproduire le défaut sur une vraie question, en diagnostiquer la cause réelle par inspection directe des données (jamais par supposition), appliquer un correctif ciblé, puis vérifier le résultat par un nouveau test réel avant de considérer le problème clos. Cette discipline, et son historique détaillé incident par incident, constitue le coeur de ce rapport.
